% !TEX program = xelatex
% ============================================================
% Data-gedreven rapport-template in rijkshuisstijl (afgeleid van oob.tex).
% Rendert de bevindingen/metadata van de generator; gebruikt de projectmacro's
% uit opmaak/packages/rapport (lakken/weglakken, plaatje, ibfindinghead, cvss,
% de kaders bestuurlijk/principe/..., severity-kleuren).
% ============================================================
\documentclass[11pt,a4paper]{article}
\usepackage[a4paper,top=20mm,bottom=22mm,left=21mm,right=21mm,headheight=15pt]{geometry}
\usepackage{fontspec}
% Lato indien beschikbaar, anders de standaardfont (portabel).
\IfFontExistsTF{Lato}{\setmainfont{Lato}\setsansfont{Lato}}{}
[% if taal == 'en' %]\usepackage[english]{babel}[% else %]\usepackage[dutch]{babel}[% endif %]
\usepackage{microtype}
\usepackage{ragged2e}
\usepackage{enumitem}
\usepackage{parskip}
\usepackage{booktabs}
\usepackage{longtable}
\usepackage{tabularx}
\usepackage{xltabular}
\usepackage{titlesec}
\usepackage{fancyhdr}
\usepackage{lastpage}
% Projectmacro's + palet + kaders + redactie (gelakt/ongelakt via \weglakken).
\usepackage{opmaak/packages/rapport}
\providecolor{grijs}{HTML}{535353}
\hypersetup{colorlinks=true,linkcolor=rijkblauw,urlcolor=rijkblauw,citecolor=rijkblauw,
  pdftitle={[[ titel ]]},pdfauthor={[[ auteur ]]}}
% Koppen in rijkshuisstijl.
\titleformat{\section}{\Large\bfseries\color{rijkblauw}}{\thesection}{0.65em}{}
\titleformat{\subsection}{\large\bfseries\color{rijkblauw}}{\thesubsection}{0.65em}{}
\titleformat{\subsubsection}{\normalsize\bfseries\color{grijs}}{\thesubsubsection}{0.65em}{}
\titlespacing*{\section}{0pt}{1.7em}{0.7em}
% Koptekst/voettekst.
\pagestyle{fancy}\fancyhf{}
\fancyhead[L]{\footnotesize\color{grijs}[[ classificatie ]]}
\fancyhead[R]{\footnotesize\color{grijs}[[ project ]]}
\fancyfoot[L]{\footnotesize\color{grijs}[[ auteur ]]}
\fancyfoot[R]{\footnotesize\color{grijs}[% if taal == 'en' %]Page[% else %]Pagina[% endif %]~\thepage{}~[% if taal == 'en' %]of[% else %]van[% endif %]~\pageref*{LastPage}}
\renewcommand{\headrulewidth}{0.3pt}\renewcommand{\footrulewidth}{0pt}
\newcolumntype{Y}{>{\RaggedRight\arraybackslash}X}

% TLP-contrast: zwarte tekst alleen op lichte achtergronden (amber/green).
[% set _cl = (classificatie or '') | upper %]
[% set _tlpblack = ('AMBER' in _cl) or ('GREEN' in _cl) or ('GROEN' in _cl) %]

\begin{document}

% ---------------- Titelblad ----------------
\thispagestyle{empty}
\IfFileExists{opmaak/dictu.png}{\begin{center}\vspace*{-0.5cm}\includegraphics[width=0.55\textwidth]{opmaak/dictu.png}\end{center}}{%
  \noindent{\color{rijkblauw}\rule{\textwidth}{3pt}}}
\vspace*{3cm}
\begin{center}
  {\color{rijkblauw}\fontsize{30}{36}\selectfont\bfseries [[ titel ]]\par}
  \vspace{0.7cm}
  {\Large\itshape\color{grijs} [[ subtitel ]]\par}
  \vspace{0.6cm}
  {\color{rijkblauw}\rule{0.42\textwidth}{1.2pt}\par}
  \vspace{2.6cm}
  {\large\renewcommand{\arraystretch}{1.5}%
  \begin{tabular}{@{}r@{\hspace{1.2em}}l@{}}
    [% if project %]\textcolor{rijkblauw}{\bfseries [% if taal == 'en' %]Client[% else %]Opdrachtgever[% endif %]} & [[ project ]] \\[% endif %]
    \textcolor{rijkblauw}{\bfseries [% if taal == 'en' %]Author[% else %]Opsteller[% endif %]} & [[ auteur ]] \\
    \textcolor{rijkblauw}{\bfseries [% if taal == 'en' %]Date[% else %]Datum[% endif %]} & [[ datum ]] \\
    [% if rapport_versie is defined and rapport_versie %]\textcolor{rijkblauw}{\bfseries [% if taal == 'en' %]Version[% else %]Versie[% endif %]} & [[ rapport_versie ]] \\[% endif %]
    [% if testscope %]\textcolor{rijkblauw}{\bfseries Scope} & [[ testscope | latex_escape ]] \\[% endif %]
  \end{tabular}\par}
  [% if is_draft %]\vspace{1.3cm}{\large\bfseries\textcolor{ored}{\fbox{\strut\,[% if taal == 'en' %]DRAFT[% else %]CONCEPT[% endif %]\,}}\par}[% endif %]
\end{center}
\vfill
\noindent\colorbox{[[ tlp_color ]]}{%
  \begin{minipage}[c][0.95cm][c]{\dimexpr\textwidth-2\fboxsep\relax}\centering
    \color{[% if _tlpblack %]black[% else %]white[% endif %]}{\large\bfseries [% if taal == 'en' %]Classification[% else %]Classificatie[% endif %]:~[[ classificatie ]]}%
  \end{minipage}}%
\newpage

% ---------------- Inhoudsopgave ----------------
\tableofcontents
\newpage

% ---------------- Managementsamenvatting ----------------
[% if management_samenvatting %]
\section*{[% if taal == 'en' %]Executive Summary[% else %]Managementsamenvatting[% endif %]}
\addcontentsline{toc}{section}{[% if taal == 'en' %]Executive Summary[% else %]Managementsamenvatting[% endif %]}
[[ management_samenvatting ]]

[% if algeheel_beeld is defined and algeheel_beeld %]\bestuurlijk{[[ algeheel_beeld ]]}[% endif %]
\newpage
[% endif %]

% ---------------- Bevindingenlijst (overzicht) ----------------
[% if bevindingenlijst_rows %]
\section{[% if taal == 'en' %]Findings overview[% else %]Bevindingenlijst[% endif %]}
{\arrayrulecolor{white}\renewcommand{\arraystretch}{1.25}%
\begin{xltabular}{\textwidth}{|c|Y|c|}
  \hline
  \rowcolor{rijkblauw}\textcolor{white}{\textbf{ID}} & \textcolor{white}{\textbf{[% if taal == 'en' %]Finding[% else %]Bevinding[% endif %]}} & \textcolor{white}{\textbf{[% if taal == 'en' %]Risk[% else %]Risico[% endif %]}} \\ \hline
  \endhead
  [% for r in bevindingenlijst_rows %]
  [[ r.id ]] & [[ r.titel | latex_escape ]] & \cellcolor{[[ r.color ]]}\textbf{\textcolor{[[ r.fg ]]}{[[ r.risk ]]}} \\ \hline
  [% endfor %]
\end{xltabular}}
[% endif %]

% ---------------- Bevindingen (detail) ----------------
\section{[% if taal == 'en' %]Findings[% else %]Bevindingen[% endif %]}
[% for severity_label, group in grouped_findings %]
[% if group %]
\subsection{[[ severity_label ]] [% if taal == 'en' %]findings[% else %]bevindingen[% endif %]}
[% for item in group %]
\ibfindinghead{[[ '%03d' | format(item.nr) ]]}%
  {[[ item.titel | latex_escape ]]}%
  {[[ item.basescore ]]}%
  {[[ (item.owasp or '—') | latex_escape ]]}%
  {[[ (('CWE-' ~ item.cwe) if item.cwe else '—') | latex_escape ]]}%
  {[[ (item.ncsc or '—') | latex_escape ]]}

\noindent [[ item.beschrijving | latex_markup ]]
[% if item.evidence_tex is defined and item.evidence_tex %]
[[ item.evidence_tex ]]
[% endif %]
[% if item.uitwerking_detail %]

\noindent [[ item.uitwerking_detail | latex_markup ]]
[% endif %]

\ibfindingsection{Impact}
\noindent [[ item.impact | latex_markup ]]

\ibfindingsection{[% if taal == 'en' %]Recommendation[% else %]Aanbeveling[% endif %]}
\noindent [[ item.aanbeveling | latex_markup ]]

[% if item.locatie %]

\noindent\textbf{[% if taal == 'en' %]Location[% else %]Locatie[% endif %]:}~\texttt{[[ item.locatie | latex_escape ]]}
[% endif %]
[% if item.cvss_vector %]

\noindent\textbf{CVSS Vector:}~\texttt{\small [[ item.cvss_vector | latex_escape ]]}
[% endif %]
[% if item.referenties_list is defined and item.referenties_list %]

\noindent\textbf{[% if taal == 'en' %]References[% else %]Referenties[% endif %]:}
\begin{itemize}
[% for ref in item.referenties_list %]
  \item \texttt{\small [[ ref | latex_markup ]]}
[% endfor %]
\end{itemize}
[% endif %]
\vspace{0.6em}
[% endfor %]
[% endif %]
[% endfor %]

% ---------------- Begrippenlijst ----------------
[% if begrippen_rows %]
\newpage
\section{[% if taal == 'en' %]Glossary[% else %]Begrippen- en afkortingenlijst[% endif %]}
{\renewcommand{\arraystretch}{1.3}%
\begin{xltabular}{\textwidth}{>{\bfseries}p{0.26\textwidth} Y}
  \toprule
  [% for r in begrippen_rows %]
  [[ r.term ]] & [[ r.definitie ]] \\
  [% endfor %]
  \bottomrule
\end{xltabular}}
[% endif %]

\end{document}
